\chapter{Floating org-mode tables}
\label{app:floating-org-tables}
\index{org-mode!floating tables}\index{sidewaystable}\index{org-export-filter-table-functions}

When you export an Org document to \LaTeX, tables sometimes overflow
the text width, especially wide comparison tables with many columns.
The function below hooks into \texttt{org-export-filter-table-functions}
and rewrites \emph{every} table as a full-width, landscape
\texttt{sidewaystable} that floats on its own page.  It handles both
captioned tables (which Org exports as
\texttt{\textbackslash begin\{table\}\dots\textbackslash begin\{tabular\}})
and uncaptioned ones (exported as \texttt{longtable}), converting both
to \texttt{tabularx} with the \texttt{X} column type for automatic
width distribution.

\medskip
\noindent
Three \LaTeX{} packages are required in the Org preamble:

\begin{minted}{elisp}
;; Add to your org file's preamble:
;; #+LATEX_HEADER: \usepackage{rotating}
;; #+LATEX_HEADER: \usepackage{tabularx}
;; #+LATEX_HEADER: \usepackage{pdflscape}
\end{minted}

\section{The export filter function}

\begin{minted}{elisp}
(defun my/org-latex-landscape-tables (text backend info)
  "Float tables as full-width sidewaystables on their own landscape pages."
  (when (org-export-derived-backend-p backend 'latex)
    (cond
     ;; Captioned tables: Org exports as \begin{table}...\begin{tabular}
     ((string-match "\\\\begin{table}" text)
      (let ((result text))
        (when (string-match "\\\\begin{tabular}{\\([^}]*\\)}" result)
          (let* ((cols (match-string 1 result))
                 (new-cols (replace-regexp-in-string "[lrc]" "X" cols)))
            (setq result
                  (replace-regexp-in-string
                   "\\\\begin{tabular}{[^}]*}"
                   (format "\\\\begin{tabularx}{\\\\linewidth}{%s}" new-cols)
                   result))
            (setq result
                  (replace-regexp-in-string
                   "\\\\end{tabular}" "\\\\end{tabularx}" result))))
        ;; table -> sidewaystable[p] (p = own page only, never inline)
        (setq result
              (replace-regexp-in-string
               "\\\\begin{table}\\[?[^]]*\\]?"
               "\\\\begin{sidewaystable}[p]" result))
        (setq result
              (replace-regexp-in-string
               "\\\\end{table}" "\\\\end{sidewaystable}" result))
        result))

     ;; Uncaptioned tables: Org exports as longtable -- wrap in landscape
     ((string-match "\\\\begin{longtable}{\\([^}]*\\)}" text)
      (let* ((cols (match-string 1 text))
             (new-cols (replace-regexp-in-string "[lrc]" "X" cols))
             (result text))
        (setq result
              (replace-regexp-in-string
               "\\\\begin{longtable}{[^}]*}"
               (format "\\\\begin{tabularx}{\\\\linewidth}{%s}" new-cols)
               result))
        (setq result
              (replace-regexp-in-string
               "\\\\end{longtable}" "\\\\end{tabularx}" result))
        (setq result
              (replace-regexp-in-string
               "\\\\endhead\\|\\\\endfoot\\|\\\\endfirsthead\\|\\\\endlastfoot"
               "" result))
        (concat "\\begin{landscape}\n" result "\n\\end{landscape}")))

     ;; Fallback: return unchanged
     (t text))))
\end{minted}

\section{Registering the filter}

Register the function so it runs automatically on every Org-to-\LaTeX{}
export:

\begin{minted}{elisp}
(add-to-list 'org-export-filter-table-functions
             #'my/org-latex-landscape-tables)
\end{minted}

\noindent
With this filter active, every table in your Org file is typeset in
landscape orientation.  Captioned tables become numbered
\texttt{sidewaystable} floats; uncaptioned ones are wrapped in a
\texttt{landscape} environment.  In both cases, \texttt{tabularx} with
\texttt{X} columns distributes the available width evenly across all
columns, eliminating horizontal overflow.

\begin{expertdepth}
The \texttt{[p]} placement specifier on \texttt{sidewaystable} forces
the float to a dedicated page.  Without it, \LaTeX{} may attempt to
place the rotated table inline, which almost never works for
full-width tables.  If you want rotated tables to appear closer to
their source text, try \texttt{[htp]} instead, but expect the
occasional misplaced float.

Note that this filter rewrites \emph{every} table unconditionally.
If you have small tables that fit within the text block, you may
prefer a more selective filter that only rotates tables above a
certain column count.  In that case, count the column separators
(\texttt{\&}) in the first row and skip the rewrite when the count is
below your threshold.
\end{expertdepth}
